\documentclass[10pt,a4paper]{amsart}
\usepackage[margin=19mm]{geometry}
\usepackage{amsmath,amssymb,mathtools}
\usepackage[scr=rsfs,cal=euler]{mathalfa}
\usepackage{xfrac}
\usepackage[unicode,hidelinks,bookmarks=false]{hyperref}
\newcommand{\ie}{i.e.\ }
\newcommand{\cf}{cf.\ }
\newcommand{\resp}{respectively\ }
\newcommand{\Subsection}{\subsection}
\providecommand{\nobreakdash}{\nobreak\hbox{-}}
\setlength{\emergencystretch}{2em}
\multlinegap0pt
\usepackage{amssymb}
\usepackage{xcolor}
\usepackage[shortlabels]{enumitem}
\usepackage{mathtools}
\usepackage{tikz}
\usepackage{float}
\newtheorem{claim}{Claim}
\usepackage{cleveref}
\crefname{claim}{Claim}{Claims}
\def\T{\mathbf{T}}
\newcommand{\m}[1]{\mathbf{\uppercase{#1}}}
\newcommand{\n}{\scriptscriptstyle{\mathsf{N}}}
\title{On 2-generated minimal Taylor algebras of size 4}
\author{Zarathustra Brady, Petar {\DJ}api\'c, Vuka\v{s}in {\DJ}inovi\'c, Petar Markovi\'c, Aleksandar Proki\'c, Veljko Tolji\'c, Vlado Uljarevi\'c}
\date{Source: arXiv:2604.05214v1 (CC BY 4.0)}
\begin{document}
\maketitle

\noindent\emph{Source and licence.} On 2-generated minimal Taylor algebras of size 4. Version arXiv:2604.05214v1 (CC BY 4.0), distributed by arXiv under the Creative Commons Attribution 4.0 International licence, http://creativecommons.org/licenses/by/4.0/, as recorded in the arXiv licence metadata for this exact version and shown on the abstract page https://arxiv.org/abs/2604.05214v1. Full licence notice: \texttt{licenses/arxiv-2604.05214-CC-BY-4.0.txt}.\par\smallskip

\noindent\textit{Editorial scope.} Counted unit: the article's claim claim (source0.tex lines 3178--3180). The article's complete original proof of it is delivered with it (source0.tex lines 3181--3885, one \texttt{proof} environment of 705 lines). No mathematics is written, repaired or re-derived: the statement and the proof are the article's own lines, in the article's own notation, and no retained line is substituted or re-flowed.  Retained uncounted as the source-local context the proof needs: lines 198--219.  Nothing was omitted from any retained span, so the package carries no omission record.  No retained line contains a citation, so the wrapper carries no bibliography and no bibliography entry.  No retained line contains a figure environment, an image-inclusion command or drawing code, so no image is delivered and the package declares no asset dependency.  Editorial formatting only, in addition: an AMS \texttt{amsart} preamble that restates the article's own declarations and macros used by the retained lines, namely the environments \texttt{claim} on separate counters, together with the standard packages those lines need.  The retained environments print in this document as Claim 1 in order of appearance, whereas the article prints the same items with the numbering of its own full document.  The complete native source of the article as distributed in the arXiv e-print for this exact version is bundled unchanged as \texttt{sources/197940/source0.tex}.
\subsubsection*{Nonconservative algebras}
The first three of the nonconservative algebras (\Cref{table:1}) have a single basic operation $g$, which is symmetric. The fourth algebra ${\m T}_4^{\n}$ is the semilattice algebra with binary operation $\wedge$ determined by the ordering $0 \leqslant 1$, $0 \leqslant 2$. The algebra ${\m T}_5^{\n}$ is actually the affine Mal'cev algebra $\mathbb{Z}^{\mathrm{aff}}_3$. We note that $\mathrm{min}^n_{a\leqslant b}$ in the table denotes the $n$-ary operation on $\{a,b\}$ which returns the minimum of the arguments with respect to the ordering $a\leqslant b$.
\begin{table}[h!]
\centering
 \begin{tabular}{|c||c|c|c|c|c|}
\hline
    Algebra & $g{\restriction_{\{0,1\}}}$ & $g{\restriction_{\{0,2\}}}$ & $g(1,1,2)$ & $g(1,2,2)$ & $g(0,1,2)$ \\
\hline\hline
     $\T^{\n}_1$  & maj& $\text{min}^3_{0\leqslant 2}$ & 1 & 0 & 0\\
\hline
    $\T^{\n}_2$  & aff & $\text{min}^3_{0\leqslant 2}$ & 0 & 1 & 1\\
\hline
    $\T^{\n}_3$  & maj & aff & 2 & 0 & 2\\
\hline
    $\T^{\n}_4$  & $\text{min}^3_{0\leqslant 1}$ & $\text{min}^3_{0\leqslant 2}$ & 0 & 0 & 0\\
\hline
    $\T^{\n}_5$  & \multicolumn{5}{|c|}{$\mathbb{Z}^{\mathrm{aff}}_3$}\\
\hline
\end{tabular}  
\caption{Nonconservative algebras.}
\label{table:1}
\end{table}

\begin{claim}
 $\{b,d\}$ cannot be an affine subalgebra; thus, it must be a majority.
\end{claim}

\begin{proof}
Assume that $\{b,d\}$ is affine. We look at the possibilities for $g(a,a,b)$ and $g(a,b,b)$. We cannot have $g(a,a,b)=b$ and $g(a,b,b)=a$ simultaneously. Hence, we have three cases left to discuss.

Suppose that $g(a,b,b)=a$ and $g(a,a,b)=d$ hold. Define the term $h$ by $$h(x,y,z)=g(x,g(x,y,y),g(x,y,z)),$$ and let $g'$ be defined by $g'(x,y,z)=g(h(x,y,z),h(y,z,x),h(z,x,y))$. 
$$
\begin{gathered}
h\left(
\begin{bmatrix}
b \\
a \\
a
\end{bmatrix}\!,\!
\begin{bmatrix}
a \\
b \\
a
\end{bmatrix}\!,\!
\begin{bmatrix}
a \\
a \\
b
\end{bmatrix}
\right)=g\left(
\begin{bmatrix}
b \\
a \\
a
\end{bmatrix}\!,\!
\begin{bmatrix}
d \\
a \\
a
\end{bmatrix}\!,\!
\begin{bmatrix}
d \\
d \\
d
\end{bmatrix}
\right)=
\begin{bmatrix}
b \\
\ast \\
\ast
\end{bmatrix},
\hspace{0.2cm}\text{ while }\\
h\left(
\begin{bmatrix}
a \\
b \\
b
\end{bmatrix}\!,\!
\begin{bmatrix}
b \\
a \\
b
\end{bmatrix}\!,\!
\begin{bmatrix}
b \\
b \\
a
\end{bmatrix}
\right)=g\left(
\begin{bmatrix}
a \\
b \\
b
\end{bmatrix}\!,\!
\begin{bmatrix}
a \\
d \\
b
\end{bmatrix}\!,\!
\begin{bmatrix}
a \\
a \\
a
\end{bmatrix}
\right)=
\begin{bmatrix}
a \\
\circ \\
a
\end{bmatrix},
\end{gathered}
$$

where $\ast\in\{b,d\}$ and $\circ\in\{a,c\}$. Thus, $g'(a,b,b)=a$ and $g'(a,a,b)=b$ for the cyclic term $g'$, a contradiction.

Let now $g(a,b,b)=c$ and $g(a,a,b)=b$. We define $w$ by $w(x,y,z)=g(x,x,g(x,y,z))$ and $g''(x,y,z)=g(w(x,y,z),w(y,z,x),w(z,x,y))$. Then
$$
\begin{gathered}
w\left(
\begin{bmatrix}
b \\
a \\
a
\end{bmatrix}\!,\!
\begin{bmatrix}
a \\
b \\
a
\end{bmatrix}\!,\!
\begin{bmatrix}
a \\
a \\
b
\end{bmatrix}
\right)=g\left(
\begin{bmatrix}
b \\
a \\
a
\end{bmatrix}\!,\!
\begin{bmatrix}
b \\
a \\
a
\end{bmatrix}\!,\!
\begin{bmatrix}
b \\
b \\
b
\end{bmatrix}
\right)=
\begin{bmatrix}
b \\
b \\
b
\end{bmatrix},
\hspace{0.2cm}\text{ while }\\
w\left(
\begin{bmatrix}
a \\
b \\
b
\end{bmatrix}\!,\!
\begin{bmatrix}
b \\
a \\
b
\end{bmatrix}\!,\!
\begin{bmatrix}
b \\
b \\
a
\end{bmatrix}
\right)=g\left(
\begin{bmatrix}
a \\
b \\
b
\end{bmatrix}\!,\!
\begin{bmatrix}
a \\
b \\
b
\end{bmatrix}\!,\!
\begin{bmatrix}
c \\
c \\
c
\end{bmatrix}
\right)=
\begin{bmatrix}
a \\
\ast \\
\ast
\end{bmatrix},
\end{gathered}
$$
where $\ast=g(b,b,c)\in \{a,c\}$. We know $g''(a,a,b)=b$. Also, If $\ast=a$, then $g''(a,b,b)=a$, which together with $g''(a,a,b)=b$ gives a contradiction. Therefore, let us assume $\ast=g(b,b,c)=c$, so $g''(a,b,b)=c$. 

Assume first that $g(b,c,c)=d$. The subcase $\mathrm{Sg}\{b,c\}=\m a$ leads to the same contradiction as in the case that assumed $g(a,b,b)=a$ and $g(a,a,b)=d$, with $a$ and $c$ swapping places. If $\mathrm{Sg}\{b,c\}=\{b,c,d\}$, this is impossible by Table~\ref{table:1}, as there is no 3-element nonconservative algebra which has both a congruence class and a factor isomorphic to $\mathbb{Z}^{\text{aff}}_2$. So it must be that $g(b,c,c)=b$ and $\mathrm{Sg}\{b,c\}=\{b,c\}$. Let $t(x,y):=g(x,y,y)$ and define the cyclic term $g'''$ by $g'''(x,y,z)=g(t(x,t(y,z)),t(y,t(z,x)),t(z,t(x,y))).$ We compute 
$$
\begin{gathered}
t\left(
\begin{bmatrix}
b \\
a \\
a
\end{bmatrix}\!,
t\left(
\begin{bmatrix}
a \\
b \\
a
\end{bmatrix}\!,\!
\begin{bmatrix}
a \\
a \\
b
\end{bmatrix}
\right)\right)=t\left(
\begin{bmatrix}
b \\
a \\
a
\end{bmatrix}\!,\!
\begin{bmatrix}
a \\
b \\
c
\end{bmatrix}
\right)=
\begin{bmatrix}
b \\
c \\
c
\end{bmatrix},
\;
t\left(
\begin{bmatrix}
a \\
b \\
c
\end{bmatrix}\!,
t\left(
\begin{bmatrix}
b \\
c \\
a
\end{bmatrix}\!,\!
\begin{bmatrix}
c \\
a \\
b
\end{bmatrix}
\right)\right)=t\left(
\begin{bmatrix}
a \\
b \\
c
\end{bmatrix}\!,\!
\begin{bmatrix}
b \\
a \\
c
\end{bmatrix}
\right)=
\begin{bmatrix}
c \\
b \\
c
\end{bmatrix},\\
t\left(
\begin{bmatrix}
a \\
c \\
b
\end{bmatrix}\!,
t\left(
\begin{bmatrix}
c \\
b \\
a
\end{bmatrix}\!,\!
\begin{bmatrix}
b \\
a \\
c
\end{bmatrix}
\right)\right)=t\left(
\begin{bmatrix}
a \\
c \\
b
\end{bmatrix}\!,\!
\begin{bmatrix}
c \\
b \\
c
\end{bmatrix}
\right)=
\begin{bmatrix}
c \\
c \\
b
\end{bmatrix}.
\end{gathered}
$$  
Hence, $g'''(a,a,b)=g'''(a,b,c)=g'''(a,c,b)=b$ and $\mathrm{Sg}\{b,c\}=\{b,c\}$ implies that $g'''(b,c,c)=b$, as well. Thus, $d\notin\mathrm{Sg}\{a,b\}$, a contradiction.

Consider now the case when $g(a,b,b)=c$ and $g(a,a,b)=d$. Let $h'(x,y,z):=g(x,x,g(x,y,y))$ and define a cyclic term $g^+$ by $$g^+(x,y,z)=g(h'(x,y,z),h'(y,z,x),h'(z,x,y)).$$ Then 
\begin{equation*}
h'\left(
\begin{bmatrix}
b \\
a \\
a
\end{bmatrix}\!,\!
\begin{bmatrix}
a \\
b \\
a
\end{bmatrix}\!,\!
\begin{bmatrix}
a \\
a \\
b
\end{bmatrix}
\right)=
g\left(
\begin{bmatrix}
b \\
a \\
a
\end{bmatrix}\!,\!
\begin{bmatrix}
b \\
a \\
a
\end{bmatrix}\!,\!
g\left(
\begin{bmatrix}
b \\
a \\
a
\end{bmatrix}\!,\!
\begin{bmatrix}
a \\
b \\
a
\end{bmatrix}\!,\!
\begin{bmatrix}
a \\
b \\
a
\end{bmatrix}
\right)
\right)=
g\left(
\begin{bmatrix}
b \\
a \\
a
\end{bmatrix}\!,\!
\begin{bmatrix}
b \\
a \\
a
\end{bmatrix}\!,\!
\begin{bmatrix}
d \\
c \\
a
\end{bmatrix}
\right)=
\begin{bmatrix}
d \\
a\\
a
\end{bmatrix}.
\end{equation*}
Now, if $g(d,a,a)=b$, then $g^+(a,a,b)=b$. If we want to avoid the immediate contradiction, we need to assume $g^+(a,b,b)=c$, but, according to the two previous paragraphs, this also leads to a contradiction. That being said, we have $g(d,a,a)=d$. The following discussion is based on what $g(a,d,d)$ is equal to. 

Suppose first that $g(a,d,d)=a$, thus $\{a,d\}$ is a subalgebra. Let $h'$ be defined as before, and we define $h''$ by $h''(x,y,z)=g(x,x,g(x,z,z))$. Then 
\begin{equation*}
h''\left(
\begin{bmatrix}
a \\
b \\
b
\end{bmatrix}\!,\!
\begin{bmatrix}
b \\
a \\
b
\end{bmatrix}\!,\!
\begin{bmatrix}
b \\
b \\
a
\end{bmatrix}
\right)=
g\left(
\begin{bmatrix}
a \\
b \\
b
\end{bmatrix}\!,\!
\begin{bmatrix}
a \\
b \\
b
\end{bmatrix}\!,\!
g\left(
\begin{bmatrix}
a \\
b \\
b
\end{bmatrix}\!,\!
\begin{bmatrix}
b \\
b \\
a
\end{bmatrix}\!,\!
\begin{bmatrix}
b \\
b \\
a
\end{bmatrix}
\right)
\right)=
g\left(
\begin{bmatrix}
a \\
b \\
b
\end{bmatrix}\!,\!
\begin{bmatrix}
a \\
b \\
b
\end{bmatrix}\!,\!
\begin{bmatrix}
c \\
b \\
d
\end{bmatrix}
\right)=
\begin{bmatrix}
a \\
b\\
d
\end{bmatrix},
\end{equation*}
\begin{equation*}
h'\left(
\begin{bmatrix}
a \\
b \\
b
\end{bmatrix}\!,\!
\begin{bmatrix}
b \\
a \\
b
\end{bmatrix}\!,\!
\begin{bmatrix}
b \\
b \\
a
\end{bmatrix}
\right)=
g\left(
\begin{bmatrix}
a \\
b \\
b
\end{bmatrix}\!,\!
\begin{bmatrix}
a \\
b \\
b
\end{bmatrix}\!,\!
g\left(
\begin{bmatrix}
a \\
b \\
b
\end{bmatrix}\!,\!
\begin{bmatrix}
b \\
a \\
b
\end{bmatrix}\!,\!
\begin{bmatrix}
b \\
a \\
b
\end{bmatrix}
\right)
\right)=
g\left(
\begin{bmatrix}
a \\
b \\
b
\end{bmatrix}\!,\!
\begin{bmatrix}
a \\
b \\
b
\end{bmatrix}\!,\!
\begin{bmatrix}
c \\
d \\
b
\end{bmatrix}
\right)=
\begin{bmatrix}
a \\
d\\
b
\end{bmatrix},
\end{equation*}
Define the term $w'$ by $$w'(x,y,z)=g(x,h'(x,y,z),h''(x,y,z)),$$ so, we have 
\begin{equation*}
w'\left(
\begin{bmatrix}
a \\
b \\
b
\end{bmatrix}\!,\!
\begin{bmatrix}
b \\
a \\
b
\end{bmatrix}\!,\!
\begin{bmatrix}
b \\
b \\
a
\end{bmatrix}
\right)=g\left(
\begin{bmatrix}
a \\
b \\
b
\end{bmatrix}\!,\!
\begin{bmatrix}
a \\
d \\
b
\end{bmatrix}\!,\!
\begin{bmatrix}
a \\
b \\
d
\end{bmatrix}
\right)=
\begin{bmatrix}
a \\
d \\
d
\end{bmatrix}.
\end{equation*}
Define the cyclic term $g^{\circ}$ by $g^{\circ}(x,y,z)=g(w'(x,y,z),w'(y,z,x),w'(z,x,y))$, thus we get $g^{\circ}(a,b,b)=a$. Now, we recall the analysis from the beginning of the proof, which implies that $\{a,b\}$ is closed for some cyclic term, a contradiction.

Therefore, we suppose that $g(a,d,d)=c$. Observe that we may assume $\mathrm{Sg}\{a,d\}=\{a,c,d\}$. Indeed, if $\mathrm{Sg}\{a,d\}=\m a$, after swapping $b$ and $d$, we find ourselves in a case where $\mathrm{Sg}\{a,b\} = \m a$ and $g(b,a,a) = b$, $g(a,b,b) = c$, which we have already ruled out. Hence, $\{a,c,d\}$ is a subalgebra and, according to \Cref{table:1}, it must be $\m t_3^{\scriptscriptstyle{\mathsf{N}}}$. Define the term $w''$ by $$w''(x,y,z)=g(x,g(x,y,y),g(x,y,z)),$$ so, we have 
 \begin{equation*}
w''\left(
\begin{bmatrix}
b \\
a \\
a
\end{bmatrix}\!,\!
\begin{bmatrix}
a \\
b \\
a
\end{bmatrix}\!,\!
\begin{bmatrix}
a \\
a \\
b
\end{bmatrix}
\right)=
g\left(
\begin{bmatrix}
b \\
a \\
a
\end{bmatrix}\!,\!
g\left(
\begin{bmatrix}
b \\
a \\
a
\end{bmatrix}\!,\!
\begin{bmatrix}
a \\
b \\
a
\end{bmatrix}\!,\!
\begin{bmatrix}
a \\
b \\
a
\end{bmatrix}
\right)\!,\!
g\left(
\begin{bmatrix}
b \\
a \\
a
\end{bmatrix}\!,\!
\begin{bmatrix}
a \\
b \\
a
\end{bmatrix}\!,\!
\begin{bmatrix}
a \\
a \\
b
\end{bmatrix}
\right)
\right)=
g\left(
\begin{bmatrix}
b \\
a \\
a
\end{bmatrix}\!,\!
\begin{bmatrix}
d \\
c \\
a
\end{bmatrix}\!,\!
\begin{bmatrix}
d \\
d \\
d
\end{bmatrix}
\right)=
\begin{bmatrix}
b \\
d\\
d
\end{bmatrix}.
\end{equation*}
Define the cyclic term $g^*$ by $g^*(x,y,z)=g(w''(x,y,z),w''(y,z,x),w''(z,x,y))$. It holds $g^*(b,a,a)=b$, but we have eliminated this case (the one with $g(a,b,b)=c$ and $g(a,a,b)=b$) before. This concludes the proof of the claim.
%Consider now the case when $g(a,b,b)=c$ and $g(a,a,b)=d$. Let $w$ and $g'$ be defined as in the previous case. Then
% $$
% w\left(
% \begin{bmatrix}
% b \\
% a \\
% a
% \end{bmatrix}\!,\!
% \begin{bmatrix}
% a \\
% b \\
% a
% \end{bmatrix}\!,\!
% \begin{bmatrix}
% a \\
% a \\
% b
% \end{bmatrix}
% \right)=g\left(
% \begin{bmatrix}
% b \\
% a \\
% a
% \end{bmatrix}\!,\!
% \begin{bmatrix}
% b \\
% a \\
% a
% \end{bmatrix}\!,\!
% \begin{bmatrix}
% d \\
% d \\
% d
% \end{bmatrix}
% \right)=
% \begin{bmatrix}
% d \\
% \circ \\
% \circ
% \end{bmatrix},
% $$where $\circ=g(a,a,d)\in\{b,d\}$, thus $g'(a,a,b)=d$. But,
% $$
% w\left(
% \begin{bmatrix}
% a \\
% b \\
% b
% \end{bmatrix}\!,\!
% \begin{bmatrix}
% b \\
% a \\
% b
% \end{bmatrix}\!,\!
% \begin{bmatrix}
% b \\
% b \\
% a
% \end{bmatrix}
% \right)=g\left(
% \begin{bmatrix}
% a \\
% b \\
% b
% \end{bmatrix}\!,\!
% \begin{bmatrix}
% a \\
% b \\
% b
% \end{bmatrix}\!,\!
% \begin{bmatrix}
% c \\
% c \\
% c
% \end{bmatrix}
% \right)=
% \begin{bmatrix}
% a \\
% \ast \\
% \ast
% \end{bmatrix},
% $$ where $\ast=g(b,b,c)\in \{a,c\}$, which, analogously to the previous case, leads to a contradiction.
\end{proof}

\end{document}
